\lecture{7}{20 March.}{}
\begin{theorem}[Dirac, 1952]
    Let \(G\) be an \(n\)-vertex graph with minimum degree at least \(\frac{n}{2}\), then \(G\) has a Hamilton cycle.    
\end{theorem}

\begin{proof}[proof of Dirac's theorem]
    We first introduce a lemma.
    \begin{lemma}
        Let \(G\) be an \(n\)-vertex graph with minimum degree \(\ge \frac{n}{2}\). Let \(P\) be a longest path in \(G\), then 
        \begin{itemize}
            \item [(1)] If \(u, v\) are the endpoint of \(P\), then \(N(u) \cup N(v) \subseteq V(P)\). 
            \item [(2)] There is a cycle \(C\) with \(V(C) = V(P)\).     
        \end{itemize}     
    \end{lemma}
    \begin{proof}
    \hfill
        \begin{itemize}
            \item [(1)] If an endpoint has a neighbor not in \(P\), we can extend \(P\) by using the edge incident to that endpoint and that neighbor, which is a contradiction.  
            \item [(2)] If \(\left\{ u, v \right\} \in E \), add it to the path to get a cycle. Otherwise, let 
            \[
                P = u = v_1 \to v_2 \to v_3 \to v_4 \to \dots \to v_i \to v_{i+1} \to \dots \to v_m = v.
            \]
            If \(\exists i\) s.t. \(\left\{ v_i, v \right\}, \left\{ u, v_{i+1} \right\} \in E  \), then take \(C\) to be 
            \[
                u \to v_{i+1} \to v_{i+2} \to \dots \to v_{m-1} \to v \to v_i \to v_{i-1} \to \dots \to v_2 \to u.
            \]  
            Such an \(i\) must exist: let \(S = \left\{ v_i : \left\{ u, v_{i+1} \right\} \in E  \right\} \), then \(\vert S \vert = \deg(u) \ge \frac{n}{2} \). Also, \(\vert N(v) \vert = \deg(v) \ge \frac{n}{2} \). Note that 
            \[
                \vert S \cup N(v) \vert \le \vert \left\{ u, v_2, \dots , v_{m-1} \right\}  \vert \le n - 1.  
            \]    
            Thus, \(S \cap N(v) \neq \varnothing\). 
        \end{itemize}
    \end{proof}

    If we have this lemma, then let \(P\) be a longest path and \(v\) is an endpoint. By (1), \(N(v) \subseteq N(P)\). Hence, \(\vert V(P) \vert \ge \frac{n}{2} + 1 \)(contain all neighbors of \(V\) and \(v\) itself). By (2), there is a cycle \(C\) with \(V(C) = V(P)\). If \(V(P) = V(G)\), then \(C\) is a Hamilton cycle. Otherwise, let \(x \in V(G) \setminus V(C)\), since \(\vert V(C) \vert = \vert V(P) \vert > \frac{n}{2}  \), and \(N(x) \ge \frac{n}{2}\) and \(\vert V(C) \cup N(x) \vert \le n \), so \(N(x) \cap V(C) \neq \varnothing \), i.e. there is a neighborhood of \(x\) in \(C\). Then we find a path longer than \(P\), contradicting our choice of \(P\).                
\end{proof}

\begin{remark}
    This minimum degree condition is sufficient but not necessary, e.g. \(G = C_n, n \ge 5\). 
\end{remark}

\begin{remark}
    The minimum degree condition is best possible. Otherwise, consider a bipartite graph \(G = (U \cup V, E)\), where \(\vert U \vert = \frac{n}{2} - 1 \) and \(\vert V \vert = \frac{n}{2} + 1 \), and for each \(u \in U\) and \(v \in V\) we have \(\left\{ u, v \right\} \in E \). Thus, this subgraph can not have a Hamilton cycle.         
\end{remark}

\begin{theorem}[Ore]
    Let \(G\) be an \(n\)-vertex graph, such that whenever \(\left\{ u, v \right\} \notin E(G) \), we have \(\deg(u) + \deg(v) \ge n\), then \(G\) has a Hamilton cycle.   
\end{theorem}